\section{Model}
\label{sec:model}

% We state our model in two parts. First, we define our environment to capture how competitors value
% each others' resources, how competition influences their willingness to exchange, and how to quantify
% the quality of an exchange. In the second section, we give a precise definition of a resource 
% exchange algorithm under a particular environment and its potential properties
% (e.g. optimality, individual rationality, and incentive compatibility).

We now present our model. In Section~\ref{sec:modelenv} we define the environment, which describes how competitors value each other's 
resources, how competition affects their willingness to trade, and how to evaluate the quality of an exchange. 
In Section~\ref{sec:modelalgo} we formally define a resource exchange algorithm within this environment and outline its desirable properties, 
such as optimality, individual rationality, and incentive compatibility.

\subsection{Environment}
\label{sec:modelenv}

\paraheadwnogap{Agents, valuations, and utilities}
There are two competitors $\competitora$ and $\competitorb$, each possessing 
a collection of $\numdatab$ and $\numdataa$ goods respectively, that the other competitor values.
Let \(
    \cbr{
        \valai
    }_{\numdatafora=0}^\numdataa,
    \big\{
        \valbj
    \big\}_{\numdataforb=0}^\numdatab
\) be a sequence of random variables, drawn independently from distributions $\priora, \priorb$, 
where $\valai$ denotes how much player $a$ values $i$ goods from $b$ and likewise for $\valbj$.
Although we assume the sequences are independent of each other, i.e. $\valai \perp \valbj$ for all $i,j$, 
each sequence could be correlated, $\valaii{i_1} \not\perp \valaii{i_2}$ for all $i_1,i_2$.
For convenience, we will take $\valaii{0} = \valbjj{0} = 0$.
Finally, while the distributions
$\priora,\priorb$,
%  $\prior$ 
are public information, neither player knows the realizations of their valuations
\(
    \cbr{
        \valai
    }_{\numdatafora=1}^\numdataa,
    \big\{
        \valbj
    \big\}_{\numdataforb=1}^\numdatab
\)
until the corresponding number of goods have been exchanged.

% \[
% \vala:\cbr{0,\ldots,\numdataa}\rightarrow\RR
% \quad
% \text{and} 
% \quad
% \valb:\cbr{0,\ldots,\numdatab}\rightarrow\RR 
% \]
% be two independent random valuation functions, capturing how much each competitor values a particular number of 
% resources received from the other.
% Equivalently, we can view these two functions as two sequence of random variables
% \(
%     \cbr{
%         \valai
%     }_{\numdatafora=0}^\numdataa,
%     \big\{
%         \valbj
%     \big\}_{\numdataforb=0}^\numdatab
% \)
% that are independent from 
% each other. For convenience, assume that 
% \(
%     \valaii{0}
%     =
%     \valbjj{0}
%     =
%     0
% \), 
% i.e. competitors derive no value from each others' resources when none are shared.

% Suppose that there are two competitors $\competitora$ and $\competitorb$, each possessing 
% a collection of $\numdatab$ and $\numdataa$ different resources
% respectively, that the other competitor values. Let 
% \[
% \vala:\cbr{0,\ldots,\numdataa}\rightarrow\RR
% \quad
% \text{and} 
% \quad
% \valb:\cbr{0,\ldots,\numdatab}\rightarrow\RR 
% \]
% be two independent random valuation functions, capturing how much each competitor values a particular number of 
% resources received from the other.
% Equivalently, we can view these two functions as two sequence of random variables
% \(
%     \cbr{
%         \valai
%     }_{\numdatafora=0}^\numdataa,
%     \big\{
%         \valbj
%     \big\}_{\numdataforb=0}^\numdatab
% \)
% that are independent from 
% each other. For convenience, assume that 
% \(
%     \valaii{0}
%     =
%     \valbjj{0}
%     =
%     0
% \), 
% i.e. competitors derive no value from each others' resources when none are shared.

Let $\setexchanges=\cbr{0,\ldots,\numdataa}\times\cbr{0,\ldots,\numdatab}$
denote the set of possible exchanges, where
\(
    \rbr{\numdatafora,\numdataforb}\in\setexchanges
\)
represents $\competitora$ receiving $\numdatafora$ of $\competitorb$'s resources and 
$\competitorb$ receiving $\numdataforb$ of $\competitora$'s resources.
Because $\competitora$ and $\competitorb$ are in competition, we assume their utilities under 
an exchange $\exchange = (i,j)\in\setexchanges$ are functions of both their own value and the value the other agent derives from the exchange.
Concretely,
let
\(
% \(
    \utilas = \utila\big(\valai,\valbj\big)
% \)
%\quad
\text{ and }
%\quad
% \(
    \utilbs = \utilb\big(\valai,\valbj\big)
% \)
\)
be the utilities for both agents under an exchange $\exchange\in\setexchanges$
where
\(
    \utila,\utilb:\RR^2\rightarrow\RR
\)
are bounded functions. 
For notational convenience,
we will denote $\utilass{0} = \utilass{(0,0)}=0$ and
$\utilbss{0} = \utilbss{(0,0)}=0$.
% we 
% assume
% \(
%     \utilass{0}
%     =
%     \utilbss{0}
%     =0
% \)
% and let the subscript $0$ refer to $(0,0)$.

This interaction between the two agents may be non-zero-sum (i.e. $\utilas \neq -\utilbs$).
Moreover, utility gains for the agents may be asymmetric;
for example, we could have
% \[
\(
    \utila(x,y)=x-\frac{1}{10}y
\)
% \quad
% \text{and}
and
% \quad
\(
    \utilb(x,y)=y-\frac{1}{3}x
\),
% \]
reflecting a scenario in which $\competitora$ is more sensitive to aiding $\competitorb$ than
vice versa.
Moreover, while the canonical use case of interest involves negative externalities due to competition (i.e. $\utila$ decreases in $\valbj$), 
our model and methods do not preclude positive externalities (i.e. $\utila$ may increase with $\valbj$).




\parahead{Welfare}
% To quantify social benefit due to an exchange, 
To quantify the quality of an exchange,
we introduce a bounded welfare function
\(
    \welfarefunc:\RR^2\rightarrow\RR
\).
We then define the welfare of an exchange $\exchange\in\setexchanges$ to be the random variable
\(
    \rewards
    =
    \welfarefunc\rbr{
        \utilas,\utilbs
    }
\).
For convenience, we will assume
\(
    \welfarefunc\rbr{0,0}=0
\).
Common examples of welfare functions  include the egalitarian welfare~\citep{mas1995microeconomic}
$\welfarefunc(x,y)=\min(x,y)$, the 
Nash welfare~\citep{nash1950bargaining} $W(x,y)=xy$, or the social welfare subject to individual rationality 
$\welfarefunc(x,y)=(x+y) 1_{\cbr{x,y\geq 0}}$.




In the definitions above, and going forward, we assume all functions are measurable unless stated otherwise.
Using the components above we define an exchange environment.

\begin{restatable}{definition}{def:exchange-environment}
    An exchange environment is a tuple
    \(
        \rbr{
            \cbr{
                \valai
            }_{\numdatafora=0}^\numdataa,
            \big\{
                \valbj
            \big\}_{\numdataforb=0}^\numdatab,
            \utila,
            \utilb,
            \welfarefunc
        }
    \),
    where
    \(
        \cbr{
            \valai
        }_{\numdatafora=0}^\numdataa
    \)
    and
    \(
        \big\{
            \valbj
        \big\}_{\numdataforb=0}^\numdatab,
    \)
    are independent sequences drawn from priors $\priora, \priorb$,
    \(
        \utila,\utilb,\welfarefunc : \RR^2\rightarrow\RR
    \)
    are bounded functions, 
    and all variables are normalized to zero as described above.
    % and 
    % \(
    %     \valaii{0},
    %     % =
    %     \valbjj{0},
    %     % =
    %     \utilass{0},
    %     % \utila(0,0),
    %     % =
    %     \utilbss{0},
    %     % \utilb(0,0),
    %     % =
    %     \welfarefunc(0,0),
    %     =0
    % \)
\end{restatable}  

\subsection{Algorithm Definiton and Properties}
\label{sec:modelalgo}

% \parahead{fe}
\parahead{Algorithm description}
Before formally defining an exchange algorithm, we first describe how an algorithm interacts with the agents.
First the algorithm will be published and agents will decide whether they wish to participate.
% in the data sharing process that the algorithm will compute. 
If both agents opt in, then the algorithm
will propose a first exchange 
\(
    \algstatett{1}
    =
    \rbr{
        \numdataforatt{1},
        \numdataforbtt{1}
    }\in\setexchanges 
\)
where $\competitora$ will receive $\numdataforatt{1}$ of $\competitorb$'s resources and 
$\competitorb$ will receive $\numdataforbtt{1}$ of $\competitora$'s resources.
After this exchange both agents truthfully report their observed values for 
each others' resources, 
\(
    \valaii{
        \numdataforatt{1}
    },
    \valbjj{
        \numdataforbtt{1}
    }
\).
Using this feedback, the algorithm will determine whether to terminate or propose another 
exchange 
\(
    \algstatett{2}
    =\rbr{
        \numdataforatt{2},
        \numdataforbtt{2}
    }\in\setexchanges 
\).
The process repeats until the algorithm terminates with a final exchange of 
\(
    \algstatett{
        \algendtime
    }
    =\rbr{
        \numdataforatt{\algendtime},
        \numdataforbtt{\algendtime}
    }\in\setexchanges 
\).
We require that exchanges are increasing\footnote{%
When we write $\rbr{x_1,y_1}\leq\rbr{x_2,y_2}$, we mean
\(x_1\leq x_2\), and \(y_1\leq y_2\).
When we write $\rbr{x_1,y_1}<\rbr{x_2,y_2}$, then at least one of the two inequalities holds strictly.
}, meaning $\algstatett{t} > \algstatett{t-1}$,
to reflect that resources cannot be unshared and that on each round the exchange algorithm does not terminate, at least one agent must receive an additional good.
We now formally define an exchange algorithm.

\begin{restatable}[Algorithm]{definition}{exchange-alg}
    \label{def:exchange-alg}
    \emph{A resource exchange algorithm is a tuple
    \(
        \rbr{
            \cbr{
                \algstatemapt
            % }_{\timeind\geq 1}, 
            }_{\timeind\geq 0}, 
            \cbr{
                \algstopmapt
            }_{\timeind\geq 0}
        }
    \)
    % of sequence of measurable functions
    where
    % The first sequence
    \begin{align*}
        \algstatemapt\rbr{
            \cbr{
                \valaii{
                    k
                }
            }_{k=0}^{
                % \numdataforatt{t-1}
                \numdataforatt{t}
            },
            \cbr{
                \valbjj{
                    k
                }
            }_{k=0}^{
                % \numdataforbtt{t-1}
                \numdataforbtt{t}
            }
        }
        =
        \rbr{
            \numdataforatt{\timeind+1},
            \numdataforbtt{\timeind+1}
        }
        =
        \algstatett{\timeind+1}
        \in\NN^2_{\geq 0}
        \quad
        \text{s.t.}
        \quad
        \algstatet
        <
        \algstatett{t+1}
        \quad
        \rbr{
            \text{here }
            \algstatett{0}
            =
            \rbr{
                \numdataforatt{0},\numdataforbtt{0}
            }
            =
            (0,0)
        }
    \end{align*}     
    % is a sequence of 
    is a sequence of 
    mappings, each from values realized for previous exchanges, to a larger exchange,
    and
    % \(
    %     \algstopmapt
    %     :
    %     % \RR^{2\timeind}
    %     \RR^{\numdataforatt{t}+1}
    %     \times
    %     \RR^{\numdataforbtt{t}+1}
    %     \rightarrow \cbr{0,1}
    % \)
    \[
        \algstopmapt
        \rbr{
            \cbr{
                \valaii{
                    k
                }
            }_{k=0}^{
                \numdataforatt{t}
            },
            \cbr{
                \valbjj{
                    k
                }
            }_{k=0}^{
                \numdataforbtt{t}
            }
        }
        \in\cbr{0,1}
    \]
    is a 
    sequence of decisions on whether to terminate proposing exchanges.
    We define
    % \[
    %     \algendtime
    %     =\inf\cbr{
    %         t\in\NN_{\geq 0} :
    %         \delta_{\timeind}\rbr{
    %             \cbr{
    %                 \valaii{k},
    %                 \valbjj{k}
    %             }_{k=0}^{\timeind-1} 
    %         }=0
    %     }
    % \] 
    \[
        \algendtime
        =\inf\cbr{
            t\in\NN_{\geq 0} :
            \delta_{\timeind}\rbr{
                \cbr{
                    \valaii{k}
                }_{k=0}^{\numdataforatt{t}}
                ,
                \cbr{
                    \valbjj{k}
                }_{k=0}^{\numdataforbtt{t}}
            }=0
        }
    \] 
    to be the time exchanging is discontinued, 
    % and require $\algendtime$ is finite. 
    and require
    \(
        \algstatett{
            \algendtime
        }
        \leq (\numdataa,\numdatab)
    \), i.e. the final exchange does not exceed the resources competitors possess.
    Let $\classalgs$ denote the set of all resource exchange algorithms. }
\end{restatable}

% Going forward, expectations $\EE$ will be with respect to the distribution of realized valuations $\vala_{i_t}, \valb_{j_t}$, welfares $\{\rewards\}_s$, and 
% decisions $\{\algstatett{t}\}_t$ due to the interaction between an algorithm $\alg$ and environment under the priors $\priora, \priorb$.
% If we wish to emphasize the dependence of $\EE$ under an algorithm $\alg$, we will write $\EE_\alg$.
% For a specified environment, the goal of the algorithm designer is to propose an algorithm 
% $\alg\in\classalgs$ to maximize the expected welfare of the final exchange,
% \(
%     \EE\sbr{
%         \rewardss{\algstatett{
%             \algendtime}
%         }
%     }
% \).

\parahead{Algorithm properties}
We now introduce desirable properties than an exchange algorithm can possess. The first such property is optimality. 
We say that an algorithm is optimal if there is no other algorithm with a higher expected welfare.

\begin{restatable}{definition}{welfare-of-alg-and-opt-alg}
    For an algorithm
    \(
        \alg
        =
        \rbr{
            \cbr{
                \algstatemapt
            }_{\timeind\geq 0}, 
            \cbr{
                \algstopmapt
            }_{\timeind\geq 0}
        }
        \in\classalgs
    \)
    we define its expected welfare to be that of the final exchange,
    \(
        \EE\sbr{
            \rewardss{
                \algstatett{\algendtime}
            }
        }
    \).
    We say that an algorithm $\alg$ is optimal if
    \begin{align*}
        % \EE_\alg\sbr{
        \EE\sbr{
            \rewardss{
                \algstatett{\algendtime}
            }
        }
        =
        \sup_{
            \alg'\in\classalgs
        }
        % \EE_{\alg'}\sbr{
        \EE\sbr{
            \rewardss{
                \algstatett{\algendtime'}'
            }
        } .
    \end{align*} 
\end{restatable}  

% In otherwords, an algorithm is optimal if there is no other algorithm with a higher expected 
% welfare. 
% The goal of a benevolent algorithm designer
% is to find an optimal $\alg\in\classalgs$ 
% to maximize the benefit of competitors collaborating
% by maximizing the expected welfare.
% Ideally the algorithm designer would like to propose an optimal exchange algorithm.
% % to maximize the benefit of competitors collaborating
% by maximizing the expected welfare.
% However, it is a priori not obvious if or when there exists an optimal resource exchange algorithm, 
% and if it does, how to go about finding it. 

Our second property for an algorithm is individual rationality and is defined as follows.
% Additionally, we would like to ensure that under the proposed algorithm is individually rational for both agents
% as defined below.

\begin{restatable}{definition}{ir-exchange-alg}
    An algorithm $\alg\in\classalgs$ is individually rational (IR) if 
    $\EE[\utilass{ \algstatett{\algendtime}
            }] \geq 0$
    and
    $\EE[\utilbss{ \algstatett{\algendtime}
            }] \geq 0$.
    % \begin{align*}
    % \(
    %     0
    %     \leq
    %     \EE\sbr{
    %         \utilass{ \algstatett{\algendtime}
    %         }
    %     },
    %     \EE\sbr{
    %         \utilbss{
    %             \algstatett{\algendtime}
    %         }
    %     }
    % \).
\end{restatable}  

%\kkcomment{State $\EE_A$ and that we wil drop goig forward.}

Thus far, we have assumed that once an algorithm is published competitors either opt in or out of the 
algorithm for the entire duration. However, in some real-world use cases, agents may exit the collaboration early. This is a potential issue when one competitor 
benefits significantly from collaboration early on and determines it would only serve their 
competitor by continuing to participate. An algorithm which has the property that competitors are 
never incentivized to leave the collaboration early (prior to the algorithm terminating) is
called incentive compatible (IC). This is made precise in the following definition.

\begin{restatable}{definition}{ic-exchange-alg}
    % We say that an algorithm 
    \label{def:ic-exchange-alg}
    Let
    \(
        \alg
        =
        \rbr{
            \cbr{
                \algstatemapt
            }_{\timeind\geq 0}, 
            \cbr{
                \algstopmapt
            }_{\timeind\geq 0}
        }
        \in\classalgs
    \) be an algorithm, 
    and let
    % is incentive compatible if for any alternate sequence of stopping decisions
    \(
        \cbr{
            \altalgstopmapt
        }_{\timeind\geq 0}
    \)
    be an alternate sequence of stopping decisions, i.e.
    \[
        \altalgstopmapt
        \rbr{
            \cbr{
                \valaii{
                    k
                }
            }_{k=0}^{
                \numdataforatt{t}
            },
            \cbr{
                \valbjj{
                    k
                }
            }_{k=0}^{
                \numdataforbtt{t}
            }
        }
        \in\cbr{0,1}
        \quad
        \text{and}
        \quad
        \altalgendtime
        =\inf\cbr{
            t\in\NN_{\geq 0} :
            \altalgstopmaptt{\timeind}\rbr{
                \cbr{
                    \valaii{k}
                }_{k=0}^{\numdataforatt{t}}
                ,
                \cbr{
                    \valbjj{k}
                }_{k=0}^{\numdataforbtt{t}}
            }=0
        } .
    \]
    We say $\alg$ is incentive compatible if for any choice of 
    \(
        \cbr{
            \altalgstopmapt
        }_{\timeind\geq 0}
    \)
    we have
    \[
        \EE\sbr{
            \utilass{
                \algstatett{
                    \algendtime\wedge\altalgendtime
                }
            }
        }
        \leq
        \EE\sbr{
            \utilass{
                \algstatett{
                    \algendtime
                }
            }
        }
        \quad
        \text{and}
        \quad
        \EE\sbr{
            \utilbss{
                \algstatett{
                    \algendtime\wedge\altalgendtime
                }
            }
        }
        \leq
        \EE\sbr{
            \utilbss{
                \algstatett{
                    \algendtime
                }
            }
        } .
    \] 
    Let $\classicalgs$ denote the set of IC exchange algorithms.
    We say that $\alg\in\classicalgs$ is IC optimal if
    \[
        \EE\sbr{
            \rewardss{
                \algstatett{\algendtime}
            }
        }
        =
        \sup_{
            \altalg\in\classicalgs
        }
        \EE\sbr{
            \rewardss{
                \algstatett{\algendtime'}'
            }
        } .
    \]
\end{restatable}  

In the definition above, any 
\(
    \cbr{
        \altalgstopmapt
    }_{\timeind\geq 0}
\)
can be viewed as a sequence of decisions a competitor makes on whether to exit the collaboration or
continue participating.
With this interpretation,
% In the definition above, 
\(
    \algendtime\wedge\altalgendtime
\)
represents the time collaboration stops, either because the algorithm has terminated on round $T$, or one of 
the participants has decided to exit early at round $T'$.
The definition above says that incentive compatible algorithms are those for which neither competitor
can increase their expected utility by trying to strategically exit the collaboration early.
As we will see, incentive compatibility is a strong requirement for an algorithm and significantly
complicates its design.